\exnumber{02}
\extitle{UX I - Ex 3 - Personas}
\exfiles{ex02/fichier Figma et PDF, nomme Name\_FirstName\_DeliverableName\_Date\_VersionNumber}
\exauthorize{Ceux nommés dans l'exercice}
\makeheaderfiles

\textbf{Instructions :}

\begin{itemize}
\item Créer un persona utilisateur à partir de plusieurs utilisateurs de la même catégorie
\end{itemize}

\textbf{Livrables :}

\begin{itemize}
\item Un persona utilisateur sur un document PDF
\item Déposez le PDF sur Github, nommé ainsi : « Name\_FirstName\_DeliverableName\_Date\_VersionNumber ». Exemple : « Doe\_John\_MidFiPrototype\_05242024\_V1 ».
\item Ajoutez un titre dans le document.
\item Soignez l'aspect global du document. Il doit être clair, simple et lisible. Vous pouvez regarder canva.com, sans surcharger le dessin de détails.
\end{itemize}

\textbf{Pistes :}

\begin{itemize}
\item Vous pouvez utiliser Figma et en extraire un document PDF.
\item Utilisez le matériel obtenu à l'exercice 1.
\item Repérer les personnes qui rencontrent le même type de problèmes et les regrouper.
\item Quand vous trouvez 2+ personnes qui rencontrent le même type de problème, vous pouvez fusionner leurs témoignages pour créer un persona. Dans ce cas, Joe, Nathan, Guillaume et Anaelle qui courent des courses de temps en temps (ils ont couru au moins une course dans leur vie, mais moins de 3 fois par an) et ont rencontré des problèmes de déshydratation.
\item Un persona est fictif : n'utilisez pas le nom d'un de vos utilisateurs, mais un prénom cohérent avec la démographie et le groupe social de vos utilisateurs. Dans ce cas, nous choisirons un homme car 3 sur 4 sont des hommes, et Pierre comme prénom pour qu'il soit différent (mais cohérent) de ceux de Joe, Nathan, Guillaume et Anaelle.
\item Pour l'âge, faites la moyenne des âges de vos utilisateurs sélectionnés. Dans ce cas, 32 (Joe a 30 ans, Nathan 39, Guillaume 30 et Anaelle 29).
\item Le persona doit être orienté sport, car c'est le contexte de cette recherche. Cependant, des informations plus larges, comme les infos de screening (métiers, localisation, etc.), sont les bienvenues pour donner plus de contexte.
\end{itemize}

\textbf{Ressources :}

\begin{itemize}
\item Rendez votre persona excellent à nouveau \url{https://uxdesign.cc/personas-e60c1c06ead1}
\item Persona, journey maps, user flow, site map, OH MY! \url{https://uxdesign.cc/personas-journey-maps-site-maps-and-user-flows-oh-my-e71d044b4bcb}
\item Framework Persona \url{https://miro.medium.com/max/1400/0*RbHfR6BHhhFIENNf}
\end{itemize}

\textbf{Citation :}

« Concentrez-vous sans relâche sur le client » Jesse Hertzberg | Ancien COO de Squarespace

